\documentclass[11pt,a4paper]{article}

% ===== Document metadata =====
\newcommand{\coursecode}{COMP2120}
\newcommand{\coursetitle}{Computer Organization}
\newcommand{\notetitle}{Lecture Notes}
\newcommand{\notedate}{May 17, 2026}

\usepackage{../styles/notebook}
\renewcommand{\notebookimagedir}{HKU COMP2120 images}

\begin{document}
\makenotestitle

\pagenumbering{roman}
\tableofcontents
\newpage
\pagenumbering{arabic}

\section{Introduction}
Computer organization is about the operational units and interconnections that realize the architectural specification. \par

\section{Evolution of Computers}
A program: a lot of instructions, each instruction has several cycles. \par
Average cycles per instruction (CPI):
\[
CPI = \frac{\sum_i CPI_i \times I_i}{I_c}, I_c = \sum_i I_i
\]
Processor time $T$:
\[
T = \frac{I_c \times CPI}{f}
\]
where $f$ is clock rate.

\section{Digital Logic}
\subsection{Boolean Algebra Operations}
AND$(A \cdot B)$, OR$(A + B)$, NOT$(\overline{A})$ \par
NAND$(\overline{A \cdot B})$ and NOR$(\overline{A + B})$: the complement of AND and OR \par
XOR$(A \oplus B)$: exactly one of the operands has the value 1 \par

\subsection{Boolean Algebra Laws}
\begin{enumerate}
    \item [1. ] Commutative law
        \[
        AB = BA, A + B = B + A
        \]

    \item [2. ] Identity elements
        \[
        1 \cdot A = A, 0 + A = A
        \]

    \item [3. ] Null law
        \[
        0 \cdot A = 0, 1 + A = 1
        \]
    
    \item [4. ] Idempotent law 
        \[
        A \cdot A = A, A + A = A
        \]
    
    \item [5. ] Inverse law
        \[
        A \cdot \overline{A} = 0, A + \overline{A} = 1
        \]
    
    \item [6. ] Associative law
        \[
        A \cdot (B \cdot C) = (A \cdot B) \cdot C
        \]
        \[
        A + (B + C) = (A + B) + C
        \]
    
    \item [7. ] Distributive law
        \[
        A (B + C) = AB + AC
        \]
        \[
        A + (BC) = (A + B)(A + C)
        \]

    \item [8.] DeMorgan's Theorem
        \[
        \overline{A \cdot B} = \overline{A} + \overline{B}, \overline{A + B} = \overline{A} \cdot \overline{B}
        \]
        \[
        A + B = \overline{\overline{A} \cdot \overline{B}}
        \]

\end{enumerate}

\subsection{Gates}
\insertcourseimage{3_1}

\subsection{Functionally Complete Sets}
$\{\text{AND}, \text{NOT}\}$, $\{\text{OR}, \text{NOT}\}$, $\{\text{NAND}\}$, and $\{\text{NOR}\}$ are functionally complete sets.

\subsection{Implementation of Boolean Function}
Sum of products (SOP): $F = XXX + XXX + \dots$ \par
Product of sums (POS): $F = (X + X + X) \cdot (X + X + X)  \cdot \dots$

\subsection{Logic Gate Implementation}
Half adder: add two single binary digits $A + B$:
\[
\left\{A, B\right\} \rightarrow \left\{C_{out}, S\right\}
\]

Full adder: add binary numbers and account for values carried in as well as out:
\[
\left\{A, B, C_{in}\right\} \rightarrow \left\{C_{out}, S\right\}
\]

\section{Number Representation}
\subsection{Radix Number Systems}
Positional number system: a string of digits
\[
(\dots a_3 a_2 a_1 a_0.a_{-1} a_{-2} a_{-3} \dots)_r
\]
Value:
\[
\cdots + a_3r^3 + a_2r^2 + a_1r^1 + a_0r^0 + a_{-1}r^{-1} + a_{-2}r^{-2} + a_{-3}r^{-3} + \cdots
\]
\[
= \sum_i (a_i \times r^i)
\]

\subsection{Sign and Magnitude Representation}
Dealing with negative numbers

\subsubsection{Excess $K$}
Also referred to as offset binary, excess code or biased representation.
$n$-bit Excess $K$ notation: range form $0-K$ to $2^n-1-K$.
\[
f(\text{bit pattern}) = \text{value} = \text{bit pattern} - K
\]
\note{bit pattern: the unsigned integer value}

\subsubsection{One's Complement}
Positive numbers: same as an unsigned bit pattern \par
Negative numbers: representing $-N$ by inverting all bits of binary representation of $N$
\[
f(\text{bit pattern}) =
\begin{cases}
    \text{bit pattern} - (2^n - 1) \text{  where leftmost bit} = 1 \\
    \text{bit pattern}
\end{cases}
\]

\subsubsection{Two's Complement}
Positive numbers: same as an unsigned bit pattern \par
Negative numbers: representing by One's complement + 1 \par
The bit pattern of $-N$ has a value of $2^n-N$. \par 

\textbf{Range extension}
\[
\begin{array}{r@{}c@{}l}
18 \phantom{0} & = & \phantom{000000000}00010010 \quad (\text{two's complement, 8 bits}) \\
18 \phantom{0} & = & \phantom{0}0000000000010010 \quad (\text{two's complement, 16 bits})
\end{array}
\]
\[
\begin{array}{r@{}c@{}l}
-18 \phantom{0} & = & \phantom{000000000}11101110 \quad (\text{two's complement, 8 bits}) \\
-18 \phantom{0} & = & \phantom{0}1111111111101110 \quad (\text{two's complement, 16 bits})
\end{array}
\]

\textbf{Subtraction rule} \par
To subtract $B$ from $A$, take the two's complement of $B$ and add it to $A$.

\subsection{Floating-point Representation}
\[
\pm (1.\text{xxxx}) \times 2^{\text{Exponent}} \text{or} \pm (0.1\text{xxxx}) \times 2^{\text{Exponent}} 
\]

How to represent a floating-point number? \par
a sign bit $\pm$ \par
\textbf{Significand}: normalized to 1.xxxx or 0.1xxxx. The first digit is always 1, which can be omitted. \par
\textbf{Exponent}: some signed number representation format (biased representation or excess $2^{m-1}-1$) \par
\insertcourseimage{4_1}

\subsection{Multiplication of integers (2's complement)}
For two $n$-bit binary integers, the product is at most $2n$ bits in length. \par
Both multiplicand and multiplier are positive numbers: perform unsigned binary integer multiplication \par
Negative multiplicand / multiplier:
\begin{enumerate}
    \item [\textbf{Step 1. }] Calculate the partial product for each bit in the multiplier \textbf{except the sign bit}.
    
    \item [\textbf{Step 2. }] Sign-extend the number to become a double precision number $2n$-bit.

    \item [\textbf{Step 3. }] \textbf{Sign bit of multiplier:} if sign bit = 0, do nothing; if sign bit = 1, take two's complement of multiplicand and sign extend. Add this to the partial sum.

    \item [\textbf{Step 4. }] Ignore carry out (only the carry out beyond $2n$ bits) during addition.

\end{enumerate}
\insertcourseimage{4_2}

\subsection{Floating-point Multiplication/Division}
\[
(\pm) m_1 \times 2^{exp_1} \times (\pm) m_2 \times 2^{exp_2} = (\pm) m_1 \times m_2 \times 2^{exp_1 + exp_2}
\]
\[
(\pm) m_1 \times 2^{exp_1} \div (\pm) m_2 \times 2^{exp_2} = (\pm) (m_1 / m_2) \times 2^{exp_1 - exp_2}
\]

\section{Instruction Execution Cycle}
\subsection{Basic information}
\textbf{Word}: the unit of organization of memory \par
\textbf{Register}: a small amount of fast storage, quickly accessible location available to a computer's processor, usually word-sized \par
\textbf{Word addressing}: addresses of memory on a computer uniquely identify words of memory\par

\subsection{Storage Inside CPU}
\textbf{PC - Program Counter}: holds the address of the instruction to be fetched next \par
\textbf{IR - Instruction Register}: contains the current instruction \par
\textbf{MAR - Memory Address Register}: specifies the address in memory for the next read or write, directly connected to the external address bus to the memory \par
\textbf{MBR - Memory Buffer Register}: contains the data returned from the external memory, or the data to be written to the external memory \par
\textbf{I/O AR - I/O Address Register}: hold the I/O address \par
\textbf{I/O BR - I/O Buffer Register}: hold the data to/from I/O \par

\subsection{Instruction Format}
Example: \par
\begin{tabular}{c|c|c|c}
    Op code & src operand 1 & src operand 2 & dst operand \\
    1 byte & 1 byte & 1 byte & 1 byte
\end{tabular}
\par
Operation: ADD R2, R3, R1 $\rightarrow$ 0x 00 02 03 01 \par

\note{\texttt{0x} means hexadecimal.} \par
LD / ST are two-word instructions; second word contains the address. \par
MOV: no src2, field set to 0x00.

\subsection{Branch Instructions}
Conditional branch depends on flags (e.g., zero flag) \par
\textbf{Conditional code (cc)} specifies branch type: \par
\quad BR (0x00): unconditional \par
\quad BZ (0x01): branch if zero \par
\quad BNZ (0x02): branch if not zero \par
Op code for branch = 0x08. Target address follows in next word. \par

\subsection{Instruction Cycle}
PC is automatically increased (e.g., +4 for 4-byte instruction) 

\subsection{Operand Fetch}
Registers: register $\rightarrow$ register file read $\rightarrow$ data to ALU \par
Memory (LD example): \par
\quad fetch address from next word (MAR $\leftarrow$ PC, MBR $\leftarrow$ mem[MAR]) \par
\quad then read again: MAR $\leftarrow$ MBR, MBR $\leftarrow$ mem[MAR] $\rightarrow$ operand obtained \par

\subsection{Execute \& Write-back}
ALU processes source operands $\rightarrow$ output to temporary register \par
\[ 
\text{Destination = }
\begin{cases}
    \text{register: ALU output $\rightarrow$ register file write} \\
    \text{memory: data $\rightarrow$, then mem[MAR] $\leftarrow$ MBR}
\end{cases}
\]

\subsection{Interrupts}
Basic cycle has \textbf{no escape} $\rightarrow$ need interrupts for: \par
\quad I/O attention (prevent data loss) \par
\quad Time-sharing (switch programs) \par
Interrupts allow stopping execution and transferring control to another program. \par

\section{Memory Hierarchy}
Core idea: achieve higher average access speed under reasonable cost \par
\insertcourseimage{6_1}

\subsection{Principle of Locality}
\quad \textbf{Temporal locality}: memory referenced in the recent past (e.g., instruction in loop) \par

\quad \textbf{Spatial locality}: memory whose addresses are near one another (e.g., arrays) \par
$\Rightarrow$ so future access can be from higher level.

\subsection{Random Access Memory (RAM)}
Traditional RAM is volatile. Must be provided with a constant power supply, otherwise the data are lost. \par
\quad \textbf{Dynamic RAM}: use transistor to store electric charges, need refreshing every few milliseconds (because charge will leak away), slower but cheaper \par
\quad \textbf{Static RAM}: use logic gates to store data (e.g., latch), no refreshing needed, faster but more expensive 

\begin{table}[htbp]
\centering
\begin{tabular}{|l|l|}
\hline
\textbf{Hierarchy} & \textbf{Techniques} \\ \hline
Registers & CMOS, word \\ \hline
Cache & SRAM, DRAM, block \\ \hline
Main memory & DRAM, virtual memory page \\ \hline
External memory & magnetic disk \\ \hline
\end{tabular}
\end{table}

\subsection{RAM, ROM, Flash}
\textbf{Semiconductor Memory Types}
\insertcourseimage{6_2}

\section{Cache Memory}
\subsection{Cache Organization}
Divide the memory and cache into equal-sized blocks. A \textbf{block} in cache is also called a cache \textbf{line}. \par
When a memory block is referenced, the entire block is moved into the cache. \par
The next time it is referenced, it will be retrieved from the cache, instead of from the main memory again. 

\subsection{Address Mapping}
The CPU gives out a 32-bit address. Each byte has an address. \par
The block size is $b$ bytes. \par
Then address $A$ is in block $\lfloor A/b \rfloor$, and it is the $(A \bmod b)$-th byte in the block.

\textbf{Fully Associative:} \par
Each main memory block can be loaded into any line of the cache. \par

\textbf{Direct-map Organization:} \par
Map each block of main memory into \textbf{only one} possible line. \par

\textbf{k-way Set Associative:} \par
A compromised method of the above, cache consists of a number of sets, each one a few lines. \par
Each cache set contains $k$ (most common: 2, max: 8) cache lines. Store each block of main memory into one certain cache set with $k$ possible line. \par

Address: \par
\insertcourseimage{7_1}

\note{k = 1 $\rightarrow$ direct-map organization \\ k = \#num of lines $\rightarrow$ fully associative} \par

\subsection{Replacement Algorithm}
FIFO - first-in-first-out; LRU - least recently used; RR - random replacement. \par
Average access time = hit time + miss rate $\times$ miss penalty

\section{External Storage}
\subsection{Magnetic Disks}
\textbf{Gap}: separates sectors \par
\textbf{Sync}: indicates the beginning of an information field and provides timing alignment \par
\textbf{Address mark}: contains data to identify the sector's number and location \par
\textbf{Data}: 512 bytes, or 4096 bytes \par
\textbf{ECC}: Error correction code \par
\insertcourseimage{8_1}

Constant angular velocity (CAV) \& Multiple zone recording (MZR) \par
\insertcourseimage{8_2}

\subsection{Disk Access}
\begin{enumerate}
    \item [1. ] Seek \par
        Move the R/W head from one cylinder to another \par
        Typical seek time: initial startup time 5~16ms, 0.2~1ms for consecutive track
    
    \item [2. ] Rotational delay \par
        Time for the required sector to rotate under the R/W head \par
        Average latency: half a revolution (180-degree rotation)
    
    \item [3. ] Data transfer time \par
        Typical transfer rate is 100MB/s, hence transfer time is $\ll$ seek + latency

\end{enumerate}

\subsection{RAID}
RAID: Redundant Array of Independent Disks \par
Files are striped across multiple disks. Redundant disk capacity is used to store parity information, which guarantees 7 days 24 hour operation in mission critical applications. 

\section{Input and Output}
\subsection{Input/Output}
External devices: not on the motherboard \par
Generic model of an I/O module:
\insertcourseimage{9_1}

Block diagram of an external device:
\insertcourseimage{9_2}
\begin{enumerate}
    \item [] Control logic: Analyse computer commands and control data transaction.
    \item [] Buffer: Store the temporal data in buffer (because the device is very slow compared with CPU), making the transaction more efficient.
    \item [] Transducer: Translate physical information (e.g. keyboard input) into digital signal.
\end{enumerate}

Block diagram of an I/O module:
\insertcourseimage{9_3}

\subsection{I/O Techniques}
\begin{enumerate}
    \item [\textbf{1. }] \textbf{Programmed I/O}
        When the processor sends a command, it must wait until the operation is finished, e.g. by repeatedly checking the status of the device.
    
    \item [\textbf{2. }] \textbf{Interrupt-driven I/O}
        The processor issues an I/O command and continues other work. The I/O module interrupts the processor when the operation is ready or complete.

    \item [\textbf{3. }] \textbf{Direct Memory Access (DMA)}
        A DMA controller transfers data directly between an I/O device and memory, with the CPU mainly setting up the transfer and handling completion.
\end{enumerate}

\begin{tabular}{c|cc}
    & No Interrupts & Use of Interrupts \\
    \hline
    I/O-to-memory transfer \\ through processor & Programmed I/O & Interrupt-driven I/O \\
    \hline
    Direct I/O-to-memory transfer & & Direct memory access (DMA) \\
\end{tabular}

\subsection{Programmed I/O}
The processor executes a program that gives it direct control of the I/O operation. \par
When the processor sends a command, it must wait until the operation is finished, e.g. by repeatedly checking the status of the device. \par
The Control and Status Register (CSR) is used to control the operation of the device and report its status. \par

\subsection{Interrupt-driven I/O}
In interrupt-driven I/O, the processor issues an I/O command, then continues to execute instructions from another process. \par
When I/O finishes, I/O module interrupts the CPU. \par
The CPU will then suspend the current program, execute the remaining I/O operations, then return to the suspended program.

\subsection{Direct Memory Access (DMA)}
DMA minimizes CPU intervention by using an intelligent device controller, e.g. the CPU tells the I/O processor to read the device and place the data in memory locations $x$ to $y$. \par
However, the CPU still uses memory $\Rightarrow$ the I/O processor may steal memory cycles from the CPU. \par
The I/O processor can request bus control so that the I/O device can transfer data through the memory bus. \par
During this period, the CPU may see an elongated memory cycle and wait for the end of that cycle. \par
\insertcourseimage{9_4}

\section{Instruction Set and Addressing}
\subsection{Elements of a Machine Instruction}
\quad Operation code, opcode \par
\quad Source operand reference \par
\quad Result operand reference \par
\quad Next instruction reference \par

\subsection{Operands}
3 operand instruction (e.g. ADD A, B, C; C $\leftarrow$ A + B) \par
2 operand instruction (e.g. ADD A, B; B $\leftarrow$ A + B) \par
1 operand instruction (e.g. ADD A; AC $\leftarrow$ AC + A; AC: accumulator) \par
0 operand instruction (e.g. ADD - pop top 2 elements from stack and ADD, then push back)

\subsection{Registers}
Two types of registers: general purpose registers \& dedicated purpose registers \par
General purpose registers - can be used freely \par
Dedicated purpose registers - such as Program Counter (PC) \par

\subsection{Types of Operations}
Instruction can be classified in to the following categories: \par
\begin{enumerate}
    \item [1. ] Data transfer - e.g. MOV(move), LD(load);
    \item [2. ] Arithmetic - e.g. ADD, SUB;
    \item [3. ] Logical - e.g. AND, OR, NOT;
    \item [4. ] Transfer of control - e.g. BR for branch, CALL;
    \item [5. ] Input/Output
    \item [6. ] Data conversion e.g. NEG, SXT(sign extend);
\end{enumerate}

Arithmetic operation - treat the operands as numbers \par
Logical operation - treat the operands as bit patterns \par

\subsection{Procedure/Function Call/Return}
A \textbf{procedure} is a self-contained computer program that is incorporated into a larger program. \par
At any point in the program the procedure may be invoked, or called. \par
One procedure has: \par
a call instruction: branches from the present location to the procedure \par
a return instruction: return from the procedure to the place from which it is called

\insertcourseimage{10_1}

\subsection{Addressing Modes}
\textbf{Addressing mode} describe the way to calculate the effective address of an operand. \par
$A = $ content of the address field in the instruction \par
$R = $ content of an address field in the instruction that refers to a register \par
$EA = $ actual address of the operand in memory \par
$(X) = $ contents of memory location $X$ or register $X$ \par

\begin{enumerate}
    \item [1. ] \textbf{Immediate}: operand is a constant in the instruction \par
        \[
        \text{Operand} = A
        \]
    \item [2. ] \textbf{Direct}: operand is in memory, address is given in the instruction \par
        \[
        EA = A
        \]
    \item [3. ] \textbf{Indirect}: operand is in memory, address of the operand is given in a register or memory location specified in the instruction \par    
        \[
        EA = (A)
        \]
    \item [4. ] \textbf{Register}: operand is in a register \par
        \[
        EA = R
        \]
    \item [5. ] \textbf{Register Indirect}: operand is in memory, address of the operand is given in a register specified in the instruction \par
        \[
        EA = (R)
        \]
    \item [6. ] \textbf{Displacement}: address of memory is given by register + an offset
        \[
        EA = A + (R)
        \]
    \item [7. ] \textbf{Stack}: operand is on top of the stack \par
        \[
        EA = \text{top of stack}
        \]
\end{enumerate}

\section{Assembly Language Programming}
\subsection{Introduction}
Each line has four fields: \par
[label]  mnemonic  operand  comment \par

\subsection{High Level v.s. Assembly Language}
High level languages:
\begin{itemize}
    \item More programmer friendly
    \item More ISA(Instruction Set Architecture) independent
    \item Each high-level statement translates to instructions in the ISA of the computer
\end{itemize}

Assembly languages:
\begin{itemize}
    \item Lower level, closer to ISA
    \item Very ISA-dependent
    \item One assembly language instruction is translated into one machine instruction by the assembler
\end{itemize}
\subsection{Assembler directives}
Assembler directives are directions to the assembler, which are not translated into machine instructions. \par
\begin{itemize}
    \item \textbf{.data}: tells the assembler to add all subsequent data to the data section
    \item \textbf{.text}: tells the assembler to add subsequent code to the text section
    \item \textbf{.globl name} makes name external to other files, for multiple files in the program
    \item \textbf{.space expression} reserves spaces, amount specified by the value of expression in bytes.
    \item \textbf{.word value1 [,value2],} put the values in successive memory locations
\end{itemize}

\subsection{Assembly Language Programming}
\begin{enumerate}
    \item [1. ] if-else structure
        Example: \par
        \begin{verbatim}
            if (a[0] >= a[1]) x = a[0];
            else x = a[1];
        \end{verbatim}
        \insertcourseimage{11_1}

    \item [2. ] repetition structure
        Example: \par
        \begin{verbatim}
            a = 0;
            for (i = 0; i < 10; i++) a+= i;
        \end{verbatim}
        \insertcourseimage{11_2}

    \item [3. ] while structure
        Example: \par
        \begin{verbatim}
            temp = 0;
            a = 1;
            while (temp < 100){
                temp += a;
                a++;
            }
        \end{verbatim}
        \insertcourseimage{11_3}
    
    \item [4.] function calls
        Use \textbf{call} and \textbf{ret} for function calls. Return address stored in the stack. \par

\end{enumerate}

\section{Operating System Support}
\subsection{OS as a User/Resource Interface}
\insertcourseimage{12_1}

\subsection{Virtual Memory}
\begin{itemize}
    \item \textbf{Physical address} - addresses that are used to actually access the memory, which can be smaller
    \item \textbf{Logical address} - addressing space of your program, e.g. if 32-bit address is used, then addressing space is from 0 to 4GB. It is the space seen by the program, but not in physical memory.
\end{itemize}

\subsection{Memory Management Unit (MMU)}
Similar to Cache memory, it is based on the principle of locality, and we use address mapping to map logical addresses (of programs) to physical addresses (of memory). \par

\subsection{Paging}
Paging is a memory management scheme that divides logical memory into fixed-size blocks called pages, and physical memory into blocks of the same size called page frames. \par

Each PTE (Page Table Entry) contains information as follows:
\begin{itemize}
    \item \textbf{V}: Valid bit, indicates if the page is in memory
    \item \textbf{P}: Protection bit, indicates the access permissions for the page
    \item \textbf{D}: Dirty bit, indicates if the page has been modified
\end{itemize}
\insertcourseimage{12_2}

Address Translation:
\insertcourseimage{12_3}

Page Replacement:
\begin{itemize}
    \item \textbf{FIFO}: First In, First Out
    \item \textbf{LRU}: Least Recently Used
\end{itemize}

\section{Processor Organization and RISC}
\subsection{Processor Organization}
The processor executes instructions via two operations:
\begin{itemize}
    \item moving data from one place to another
    \item performing data transformation/processing through ALU
\end{itemize}

Data movement can be described as follows: e.g. Instruction fetch:
\begin{align*}
\text{MAR} &\leftarrow \text{PC} \\
\text{PC} &\leftarrow \text{PC} + 4 \\
\text{IR} &\leftarrow \text{mem}[\text{MAR}]
\end{align*}

There are 3 stages in an instruction cycle:
\begin{enumerate}
    \item [1. ]Fetch Instruction (FI)
    \item [2. ]Decode Instruction (DI)
    \item [3. ]Instruction Execution
        \begin{itemize}
            \item Calculate Operand Addresses (CO)
            \item Fetch Operands from registers/memory (FO)
            \item Execute Instruction (EI)
            \item Write Operand (WO)
        \end{itemize}
\end{enumerate}

\subsection{Instruction Pipeline}
\insertcourseimage{13_1}
\note{Not all instructions require all stages.}

\textbf{Pipeline hazards} are the situations that prevent the next instruction from entering the execution.
\begin{itemize}
    \item Resource hazard
    \item Data hazard
    \item Control hazard
\end{itemize}

\subsection{Processor Performance}
\[
\text{Execution Time} = \text{Instruction Count} \times \text{CPI} \times \text{Clock Cycle Time}
\]
where CPI = Clock per Instruction

\subsection{Characteristics of Modern Processors}
The concept of Reduced Instruction Set Computer (RISC) is a design philosophy that emphasizes simplicity and efficiency in processor design.
\begin{itemize}
    \item All instructions are register-register type except LOAD and STORE.
    \item Fixed length and simple, fixed format instructions.
    \item Relatively few operations and addressing modes.
    \item Use of hardwired rather than microprogrammed control.
\end{itemize}
\end{document}

